\section{2. Folgen}
\Def \key{Folge} $(a_n)_{n \in \NN_0}$, $(a_n)_{n\ge 0}$ oder $(a_n)_{n=0}^\infty$. 
Abbildung $f : \NN_0 \to \RR$ (oder $f : \NN \to \RR$). \key{Bilder} $a(n) = a_n$
sind \key{Glieder der Folge}

\Def Folge \key{direkt} beschreiben: $a_n = f(n)$;
oder \key{rekursiv}: $a_n = g(a_{n-1}, a_{n-2}, \ldots)$.
$f$ und $g$ heissen \key{Bildungsgesetz}.

\Def Folge \key{konvergent}, falls \key{Grenzwert} $L \in \RR$ mit
$\forall \varepsilon > 0 \, \exists N > 0 \, \forall n > N : |a_n - L| < \varepsilon$
existiert, sonst \key{divergent}.
Man schreibt $\lim\limits_{n \to \infty} a_n = L$

\Satz Konvergente Folge besitzt genau einen eindeutigen Grenzwert

\Def Falls $\forall M > 0 \, \exists N > 0 \, \forall n > N : a_n > M$ gilt,
Folge divergiert \key{gegen unendlich}: $\lim\limits_{n\to\infty}a_n = \infty$

Falls $\forall M > 0 \, \exists N > 0 \, \forall n > N : a_n < -M$ gilt,
Folge divergiert \key{gegen minus unendlich}: $\lim\limits_{n\to\infty} a_n = -\infty$

\Def Sei $(a_n)_{n \in \NN_0}$ Folge in $\RR$.
\key{Teilfolge} ist Folge der form $(a_{n_k})_{k \in \NN_0}$
wobei $(n_k)_{k \in \NN_0}$ Folge in $\NN_0$ mit $n_{k+1} > n_k$

\Satz $(a_n)_{n \in \NN_0}$ konvergent gegen $L \in \RR$. Auch jede Teilfolge $\to L$

\Def \key{Häufungspunkt} $A$, falls
$\forall \varepsilon > 0 \, \forall N \in \NN_0 \, \exists n\ge N : |a_n - A| < \varepsilon$ \\
Äquiv.: $a_n \in (A-\varepsilon, A+\varepsilon)$ gilt für unendlich viele $n$

\Satz $A$ Häufungspunkt $\iff$ Teilfolge mit $\lim\limits_{k \to \infty} a_{n_k} = A$ existiert

\Kor Konvergente Folge hat genau einen Häufungspunkt: $L$

\Kor Sei $A$ Häufungspunkt, unendlich viele $a_n \in (A-\varepsilon, A+\varepsilon)$

\Satz Sei $\lim\limits_{n\to\infty} a_n = K \in \RR$,
$\lim\limits_{n\to\infty} b_n = L \in \RR$ und $C \in \RR$:
\vspace{-0.6em}
\begin{itemize}
    \item $\lim\limits_{n\to\infty}(a_n + b_n) = K + L$
    \item $\lim\limits_{n\to\infty}(a_n - b_n) = K - L$
    \item $\lim\limits_{n\to\infty}(a_n \cdot b_n) = K \cdot L$
    \item $\lim\limits_{n\to\infty}(C \cdot a_n) = C \cdot K$
    \item Falls $L\ne 0$, $b_n\ne 0$: $\lim\limits_{n\to\infty}\frac{a_n}{b_n} = \frac{K}{L}$
    \item Falls $K < L$, es gibt $N \in \NN$ mit $a_n < b_n$ für alle $n\ge N$
    \item Falls es $N \in \NN$ so dass $a_n\le b_n$ für alle $n\ge N$, dann gilt $K\le L$
\end{itemize}

\Satz \key{Wichtige Grenzwerte}
\vspace{-0.6em}
\begin{itemize}
    \item $\lim\limits_{n\to\infty} \frac{\log n}{n} = \lim\limits_{n\to\infty} \frac{\ln n}{n} =0$
    \item $\lim\limits_{n\to\infty} n^\frac{1}{n} = \sqrt[n]{n} = 1$
    \item $\lim\limits_{n\to\infty} x^\frac{1}{n} = \sqrt[n]{x} = 1$, $x > 0$
    \item $\forall x \in \RR : \lim\limits_{n\to\infty} (1 + \frac{x}{n})^n = e^x$
    \item $\forall x \in \RR : \lim\limits_{n\to\infty} \frac{x^n}{n!} = 0$
\end{itemize}

\Satz $\star$ Für grosse $x \in \RR$ resp. $n \in \NN$, $a, p \in \RR$, $a > 1$:
$$\ln(\ln x) \ll \ln x \ll x^p \ll a^x \ll n! \ll n^n$$

\Satz \key{Sandwich-Theorem}:
Seien $(a_n)_{n\in\NN_0}$ und $(c_n)_{n\in\NN_0}$ mit gleichem Grenzwert $L$ und
$(b_n)_{n\in\NN_0}$ mit $\forall n \in \NN_0 : a_n\le b_n\le c_n$. $\lim\limits_{n\to\infty} b_n = L$.

\Def \key{Nach unten}/\key{oben beschränkt}, falls $M = \{a_n \, | \, n \in \NN_0\}$
nach unten/oben beschränkt. Unten und oben beschränkt: \key{beschränkte Folge}

\Def \key{monoton fallend}/\key{wachsend},
falls $\forall n,m \in \NN_0 : m > n \implies a_m\le a_n$ (resp. $\geq$).
\key{Streng monoton fallend}/\key{wachsend} falls $<$ oder $>$

\Satz Jede beschränkte und monotone Folge konvergiert.

\Satz Jede konvergente Folge ist beschränkt

\Satz Eigenschaften monotoner Folgen
\vspace{-0.6em}
\begin{itemize}
    \item Eine monotone Folge konvergiert $\iff$ sie ist beschränkt
    \item Falls monoton wachsend, gilt $\lim\limits_{n\to\infty}a_n = \sup\{a_n \, | \, n \in \NN_0\}$
    \item Falls monoton fallend, gilt $\lim\limits_{n\to\infty}a_n = \inf\{a_n \, | \, n \in \NN_0\}$
\end{itemize}

\Def \key{Limes superior}: $\limsup\limits_{n\to\infty} a_n = \varlimsup\limits_{n\to\infty} a_n = \lim\limits_{n\to\infty} \sup\{a_k \, | \, k\ge n\}$
\key{Limes inferior}: $\liminf\limits_{n\to\infty} a_n = \varliminf\limits_{n \to \infty} a_n = \lim\limits_{n\to\infty} \inf\{a_k \, | \, k\ge n\}$

\Satz Es gelten die folgenden Tatsachen:
\vspace{-0.6em}
\begin{itemize}
    \item $\limsup\limits_{n\to\infty} a_n = \inf\limits_{n\in\NN_0}(\sup\{a_k \, | \, k\ge n\})$
    \item $\liminf\limits_{n\to\infty} a_n = \sup\limits_{n\in\NN_0}(\inf\{a_k \, | \, k\ge n\})$
    \item konvergente Folge: $\lim\limits_{n\to\infty} a_n = \limsup\limits_{n\to\infty} a_n = \liminf\limits_{n\to\infty} a_n$
    \item beschr. konvergiert gdw $\lim\limits_{n\to\infty}a_n = \limsup_{n\to\infty} a_n = \liminf\limits_{n\to\infty} a_n$
\end{itemize}

\Satz Beschränkt, $\limsup\limits_{n\to\infty}a_n = A$. Dann $A$ Häufungspunkt,  $\forall\varepsilon > 0$:
\vspace{-0.6em}
\begin{itemize}
    \item nur endlich viele $a_n$ mit $a_n > A + \varepsilon$
    \item unendlich viele $a_n$ mit $A - \varepsilon < a_n < A + \varepsilon$
\end{itemize}

\Kor Jede beschr. Folge hat Häufungspunkt und konv. Teilfolge

\Def \key{Cauchy-Folge}: $\forall \varepsilon > 0 \, \exists N \in \NN \, \forall m, n\ge N : |a_n - a_m| < \varepsilon$

\Satz Jede Cauchy-Folge ist beschränkt

\Satz Folge konvergiert $\iff$ ist Cauchy-Folge
